% ===========================================================================
%  PARKED --- not compiled, not deleted.
%
%  Source: CH10/document MAIN/sections/04_public_good_accumulation_a_model_and_potential_relevance_to_.tex,
%  lines 41-65: the closing subsection "Notes --- public good in a closed
%  group", including the Merchant of Venice / historical lending-practices
%  passage and the commented-out socialism--libertarianism paragraph.
%
%  Why parked: confirmed cut in CH9_DECISIONS.md section 4.  Both the
%  Shakespeare passage and the commented-out paragraph are named there
%  explicitly as cut.  Nothing from this file is carried into the public-good
%  section.
%
%  Note that CH10 retains its own copy of the whole file; removing the
%  migrated sections from CH10 is not this rewrite's job.
% ===========================================================================

\subsection*{Notes -- public good in a closed group}

What if there is some replication advantage to a cohort of Y. pestis bacteria living within the enclosure of a macrophage; and what if there is something like a public good produced by the bacteria, and confined within the boundary of the host cellular membrane -- this underpinning a benefit of within group (macrophage) cooperative co-existence. \par 

\vspace{10pt}
\noindent
Something that comes to mind is,  as in other cases of cultural practice,  in Jewish communities it has historically been virtuous to cooperate within the community whilst facing outwards with more competitiveness.  In Shakespeare's "The merchant of venice" it is revealed that at the time (16h CE) it was forbidden for Jews to lend money with interest within the community,  whilst being acceptable to do so with Christians. \par

\vspace{10pt}
\noindent
Just like the individual is rewarded -- paid by gains of cooperative advantage - for being effective to the purpose and success of the private group; so the private group can be rewarded for benefiting circles of wider scope.\par 
And going in the other direction, the individual within the individual is crucial in its avarice when chasing personal enterprise. \par 
Surely these things don't need to be contradictory. 



%My feeling is that somewhere in all this is a helathy middle ground between the two extremes of a particular ideological spectrum.  Crudely,  that between the socialist and the individualistic libertarian capitalist.  On the one hand we have everyone driven to povery by the abolition of accute progress; on the other hand, we have the pathology of extreme inequality.  Something in the middle might see a framework of nested groups internally cooperating and externally competing. 



%%%%%%%%%%%%%%%%%%%%%%%%%%%%%%%%%%%%%%%%%%%%%%%%%%%%%%%%%%%%%%%%%%%%%%%%%%%%%%%%%%%%%%%%%%



%%%%%%%%%%%%%%%%%%%%%%%%%%%%%%%%%%%%%%%%%%%%%%%%%%%%%%%%%%%%%%%%%%%%%%%%%%%%%%%%%%%%%%%%%%
